\documentclass{../templateWriting/cdwriting}

\newcommand{\authorinfo}{THE FORUM CENTER - NGUYỄN HOÀNG HUY}
\newcommand{\documenttitle}{TÀI LIỆU CHUYÊN ĐỀ WRITING}
\newcommand{\collectiontitle}{Level 2 — W4 (Jan 19-25)}

\begin{document}

\thispagestyle{firstpage}
\makeworkshopheader{\authorinfo}{\documenttitle}{\collectiontitle}
\pagestyle{otherpages}

\workshopprompt{Cycling is a healthier and more environmentally friendly form of transport. Nevertheless, cycling is getting less popular. What are the reasons of this trend? What can be done to make cycling more popular?}

\cdgreenheading{I. Phân tích đề bài \& Định hướng}
\begin{itemize}
\item \cdkv{Chủ đề}{Giao thông (Transport) \& Sức khỏe/Môi trường (Health/Environment).}
\item \cdkv{Vấn đề}{Đi xe đạp tốt nhưng ngày càng ít người đi.}
\item \cdblue{\textbf{Yêu cầu:}}
\begin{enumerate}
\item Tìm nguyên nhân (Causes).
\item Đề xuất giải pháp (Solutions).
\end{enumerate}
\item \cdkv{Mục tiêu}{Tránh liệt kê ý vụn vặt. Cần tập trung vào 2 nguyên nhân lớn và 2 giải pháp tương ứng, phát triển ý sâu hơn bằng giải thích và ví dụ.}
\end{itemize}

II. Brainstorming – Lên ý tưởng

\cdgreenheading{A. Nguyên nhân chính \textit{(Causes)}}

Tại sao người ta ngại đi xe đạp?

\cdstep{1. Lo ngại về an toàn và cơ sở hạ tầng \textit{(Safety Concerns and Poor Infrastructure)}}
\begin{itemize}
\item \cdkv{Thiếu làn đường riêng}{Hầu hết các thành phố lớn không có làn đường dành riêng cho xe đạp (dedicated cycling lanes). Người đi xe đạp phải đi chung với ô tô, xe máy tốc độ cao.}
\item \cdkv{Nguy cơ tai nạn}{Cảm giác không an toàn, sợ va chạm giao thông khiến người dân (đặc biệt là phụ huynh và người già) e ngại.}
\end{itemize}

\cdstep{2. Sự bất tiện và lối sống hiện đại \textit{(Inconvenience and Modern Lifestyle)}}
\begin{itemize}
\item \cdkv{Khoảng cách và thời gian}{Các thành phố mở rộng, khoảng cách từ nhà đến chỗ làm xa. Đi xe đạp tốn nhiều thời gian và sức lực hơn so với xe máy hoặc ô tô.}
\item \cdkv{Yếu tố thời tiết \& trang phục}{Thời tiết khắc nghiệt (quá nóng hoặc mưa) làm người đi xe đạp bị ướt hoặc đổ mồ hôi, gây bất tiện khi đến văn phòng làm việc.}
\end{itemize}

\cdgreenheading{B. Giải pháp \textit{(Solutions)}}

Làm sao để khuyến khích họ quay lại với xe đạp?

\cdstep{1. Cải thiện cơ sở hạ tầng \textit{(Improving Infrastructure)}}
\begin{itemize}
\item \cdkv{Xây dựng làn đường riêng}{Chính phủ cần đầu tư xây dựng hệ thống đường dành riêng cho xe đạp để đảm bảo an toàn.}
\item \cdkv{Hệ thống xe đạp công cộng}{Phát triển các trạm thuê xe đạp giá rẻ (public bike-sharing schemes) kết nối với trạm xe buýt/tàu điện để thuận tiện cho việc di chuyển ngắn.}
\end{itemize}

\cdstep{2. Chính sách khuyến khích và thay đổi nhận thức \textit{(Incentives and Awareness)}}
\begin{itemize}
\item \cdkv{Tiện ích tại nơi làm việc}{Khuyến khích các công ty lắp đặt phòng tắm (showers) hoặc chỗ để xe an toàn cho nhân viên.}
\item \cdkv{Tuyên truyền lợi ích}{Các chiến dịch truyền thông nhấn mạnh vào lợi ích sức khỏe (giảm béo phì) và bảo vệ môi trường (giảm khí thải).}
\end{itemize}

III. Hướng dẫn viết chi tiết

\cdstep{1. Mở bài \textit{(Introduction)}}
\begin{itemize}
\item \cdkv{Nhiệm vụ}{Paraphrase đề bài + Nêu Thesis statement (bài viết sẽ bàn về nguyên nhân và giải pháp).}
\item \cdkv{Cấu trúc ngữ pháp cần dùng}{Mệnh đề nhượng bộ (Although/Despite...) để tạo sự tương phản.}
\end{itemize}

\cdblue{\textbf{Ví dụ:}}

It is undeniable that cycling is an eco-friendly and healthy mode of transport. However, its popularity has significantly declined in recent years. This essay will analyze the primary reasons for this trend, such as safety concerns and inconvenience, before proposing feasible solutions to encourage more people to cycle.

\cdkv{(Bản dịch ý chính}{Không thể phủ nhận rằng đạp xe là phương thức giao thông thân thiện môi trường và tốt cho sức khỏe. Tuy nhiên, sự phổ biến của nó đã giảm đáng kể... Bài viết này sẽ phân tích các lý do chính... trước khi đề xuất các giải pháp khả thi...)}

\cdstep{2. Thân bài \textit{(Body Paragraphs)}}

\cdkv{Đoạn Thân bài 1}{Nguyên nhân (Causes)}
\begin{itemize}
\item \cdkv{Câu chủ đề (Topic Sentence)}{Nêu rõ có 2 lý do chính: an toàn và sự tiện lợi.}
\item \cdblue{\textbf{Triển khai ý (Logic):}}
\begin{itemize}
\item \cdkv{Ý 1 (An toàn)}{Đường xá đông đúc \$\textbackslash\{\}rightarrow\$ Nguy hiểm \$\textbackslash\{\}rightarrow\$ Mọi người sợ.}
\item \cdkv{Ý 2 (Tiện lợi)}{Cuộc sống bận rộn \$\textbackslash\{\}rightarrow\$ Cần di chuyển nhanh \$\textbackslash\{\}rightarrow\$ Xe đạp quá chậm và mệt.}
\end{itemize}
\end{itemize}

\cdblue{\textbf{Ví dụ:}}

The decline in cycling can be attributed to two main factors. First and foremost, safety is a major concern. In many modern cities, there is a lack of dedicated cycling lanes, forcing cyclists to share the road with fast-moving vehicles. This makes people feel vulnerable and afraid of accidents. Furthermore, cycling is often seen as impractical for modern life. Since people often live far from their workplaces, riding a bike can be time-consuming and physically demanding, especially in bad weather conditions.

\cdkv{(Bản dịch}{Sự sụt giảm việc đạp xe có thể quy cho 2 yếu tố chính. Đầu tiên và quan trọng nhất, an toàn là mối lo lớn... thiếu làn đường riêng... khiến mọi người cảm thấy dễ bị tổn thương. Hơn nữa, đạp xe thường bị xem là thiếu thực tế... tốn thời gian và tốn sức...)}
\begin{itemize}
\item \cdkv{Lưu ý ngữ pháp}{Sử dụng các từ nối First and foremost, Furthermore để tăng tính liên kết (Cohesion).}
\end{itemize}

\cdkv{Đoạn Thân bài 2}{Giải pháp (Solutions)}
\begin{itemize}
\item \cdkv{Câu chủ đề (Topic Sentence)}{Chính phủ và doanh nghiệp cần hành động để giải quyết vấn đề này.}
\item \cdblue{\textbf{Triển khai ý (Logic):}}
\begin{itemize}
\item \cdkv{Giải pháp 1 (Hạ tầng)}{Nhà nước xây đường riêng \$\textbackslash\{\}rightarrow\$ An toàn hơn \$\textbackslash\{\}rightarrow\$ Dân dám đi.}
\item \cdkv{Giải pháp 2 (Tiện ích)}{Công ty hỗ trợ phòng tắm/chỗ để xe \$\textbackslash\{\}rightarrow\$ Giải quyết vấn đề mồ hôi/bất tiện.}
\end{itemize}
\end{itemize}

\cdblue{\textbf{Ví dụ:}}

To reverse this trend, several measures should be implemented. The most effective solution is for the government to invest in better infrastructure. By constructing safe lanes solely for bicycles, authorities can ensure the safety of riders, which would encourage more citizens to choose this green transport. In addition, employers can play a role by providing facilities such as showers and secure parking spaces. This would address the issue of personal hygiene and convenience for employees who cycle to work.

\cdkv{(Bản dịch}{Để đảo ngược xu hướng này... Giải pháp hiệu quả nhất là chính phủ đầu tư vào hạ tầng... Bằng cách xây làn đường an toàn... Ngoài ra, người sử dụng lao động có thể đóng góp bằng cách cung cấp cơ sở vật chất như vòi hoa sen... giải quyết vấn đề vệ sinh cá nhân...)}
\begin{itemize}
\item \cdkv{Lưu ý ngữ pháp}{Sử dụng cấu trúc By + V-ing (Bằng cách làm gì đó) để chỉ cách thức thực hiện giải pháp.}
\end{itemize}

\cdstep{3. Kết bài \textit{(Conclusion)}}
\begin{itemize}
\item \cdkv{Nhiệm vụ}{Tóm tắt lại nguyên nhân (an toàn, bất tiện) và giải pháp (đầu tư hạ tầng, hỗ trợ từ công ty).}
\item \cdkv{Lưu ý}{Không đưa thêm ý mới.}
\end{itemize}

\cdblue{\textbf{Ví dụ:}}

In conclusion, the main reasons why fewer people ride bikes are the fear of traffic accidents and the inconvenience of long distances. However, by improving road safety and offering better facilities at workplaces, society can promote cycling as a primary means of transport again.

\cdkv{(Bản dịch}{Tóm lại, lý do chính... là nỗi sợ tai nạn và sự bất tiện... Tuy nhiên, bằng cách cải thiện an toàn đường bộ và cung cấp cơ sở vật chất tốt hơn... xã hội có thể thúc đẩy đạp xe trở lại...)}

IV. Từ vựng nâng cao

Để đạt điểm cao hơn, hãy thay thế các từ cơ bản bằng các từ vựng dưới đây:

Từ/Cụm từ (Word/Phrase)

Nghĩa (Tiếng Việt)

Ví dụ câu (Example)

Eco-friendly (adj)

Thân thiện với môi trường

Cycling is an eco-friendly alternative to driving cars.

Mode of transport (n)

Phương thức giao thông

We need to promote sustainable modes of transport.

Dedicated cycling lanes(n)

Làn đường dành riêng cho xe đạp

The city should build more dedicated cycling lanes to ensure safety.

Vulnerable (adj)

Dễ bị tổn thương/nguy hiểm

Cyclists are vulnerable when sharing the road with trucks.

Impractical (adj)

Thiếu thực tế/Không khả thi

Riding a bike for 20km is impractical for many office workers.

Physically demanding(adj)

Đòi hỏi nhiều thể lực

Cycling uphill can be physically demanding.

Invest in infrastructure(v)

Đầu tư vào cơ sở hạ tầng

The government must invest in infrastructure to support cyclists.

Implement measures (v)

Thi hành các biện pháp

Strict measures should be implemented to reduce traffic congestion.

Sedentary lifestyle (n)

Lối sống thụ động (ít vận động)

Cycling helps combat the sedentary lifestyle of modern workers.

Traffic congestion (n)

Tắc nghẽn giao thông

More bikes mean less traffic congestion.

Commute (v/n)

Việc đi lại (từ nhà đến chỗ làm)

His daily commute takes 30 minutes by bike.

\end{document}
